\documentclass[11pt]{article}
\usepackage[utf8]{inputenc}
\usepackage[spanish]{babel}
\usepackage{geometry}
\usepackage{amsmath}
\usepackage{booktabs}
\usepackage{xcolor}
\usepackage{mdframed}
\usepackage{hyperref}

\geometry{a4paper, margin=1in}
\hypersetup{
    colorlinks=true,
    linkcolor=blue,
    urlcolor=cyan,
}

\title{Guía Rápida: RF Switches para Monitoreo SDR}
\author{Ingeniería de Radiofrecuencia}
\date{\today}

\newmdenv[
  backgroundcolor=red!8,
  linecolor=red!60,
  linewidth=1pt,
  roundcorner=4pt,
  innertopmargin=6pt,
  innerbottommargin=6pt
]{advertencia}

\newmdenv[
  backgroundcolor=yellow!10,
  linecolor=yellow!60!black,
  linewidth=1pt,
  roundcorner=4pt,
  innertopmargin=6pt,
  innerbottommargin=6pt
]{pendiente}

\begin{document}

\maketitle
\tableofcontents
\newpage

% =========================================================
\section{Criterio de Selección}
% =========================================================

Esta guía resume qué ofrece cada switch para un sistema SDR. La idea no es documentar cada parámetro de laboratorio, sino decidir rápido qué módulo conviene usar según el montaje.

\begin{table}[h]
    \centering
    \small
    \begin{tabular}{p{0.24\textwidth}p{0.32\textwidth}p{0.34\textwidth}}
        \toprule
        \textbf{Switch} & \textbf{Útil para} & \textbf{Limitación principal} \\
        \midrule
        Opera Cake & Automatizar antenas/filtros con HackRF One. & Depende del ecosistema HackRF y de su header. \\
        Módulo SP4T F2914 & Conmutación RF general de alta frecuencia y mejor desempeño eléctrico. & Requiere control lógico externo y cableado propio. \\
        \bottomrule
    \end{tabular}
\end{table}

% =========================================================
\section{Opera Cake}
% =========================================================

Switch RF diseñado como expansión directa para \textbf{HackRF One}. Es la opción más simple si el sistema ya usa HackRF y se quiere seleccionar antenas o filtros desde software.

\begin{table}[h]
    \centering
    \small
    \begin{tabular}{p{0.30\textwidth}p{0.58\textwidth}}
        \toprule
        \textbf{Parámetro} & \textbf{Resumen utilizable} \\
        \midrule
        Tipo & Matriz con 2 puertos primarios A/B y 8 puertos secundarios. \\
        Rango nominal & 1 MHz a 6 GHz. \\
        Control & \texttt{hackrf\_operacake}. \\
        Modos útiles & Manual, por frecuencia y por tiempo. \\
        Mejor uso & Elegir automáticamente antenas o filtros según banda de monitoreo. \\
        Integración & Se monta sobre el HackRF One; no requiere controlador externo adicional. \\
        \bottomrule
    \end{tabular}
\end{table}

\textbf{Cuándo elegirlo:} cuando se necesita una solución rápida, integrada al HackRF, para alternar varias antenas o bancos de filtros en monitoreo espectral.

\textbf{Qué verificar:} pérdida de inserción, aislamiento, potencia máxima y continuidad DC de cada ruta antes de usarlo con TX, Bias-T o señales fuertes.

% =========================================================
\section{Módulo SP4T F2914}
% =========================================================

El módulo del enlace de Amazon corresponde a un switch RF tipo \textbf{SP4T} basado en la familia \textbf{F2914}. Es una opción más genérica que Opera Cake: no depende de HackRF y puede integrarse con GPIO, microcontrolador, FPGA o una placa de control.

\begin{table}[h]
    \centering
    \small
    \begin{tabular}{p{0.30\textwidth}p{0.58\textwidth}}
        \toprule
        \textbf{Parámetro} & \textbf{Resumen utilizable} \\
        \midrule
        Tipo & SP4T absorptivo, 50 $\Omega$. \\
        Rango RF & 50 MHz a 8 GHz. \\
        Puertos & 1 común hacia 4 rutas RF. \\
        Control & Tres líneas lógicas compatibles con 3.3 V o 1.8 V. \\
        Alimentación & Fuente positiva simple de 2.7 V a 5.5 V. \\
        Desempeño típico & Baja pérdida, alto aislamiento y buena linealidad. \\
        Potencia & Hasta 33 dBm CW en ruta seleccionada; 27 dBm en rutas terminadas. \\
        Mejor uso & Bancos de filtros, selección de LNAs, rutas de prueba o matrices RF compactas. \\
        \bottomrule
    \end{tabular}
\end{table}

\textbf{Cuándo elegirlo:} cuando se necesita un switch RF independiente del SDR, con mejor rango superior que Opera Cake y control digital simple.

\textbf{Qué verificar:} pinout del módulo vendido, polaridad de alimentación, tabla de control lógico y disipación si se trabaja cerca del límite de potencia.

% =========================================================
\section{Recomendación Rápida}
% =========================================================

\begin{itemize}
    \item \textbf{Para HackRF One y automatización simple:} usar Opera Cake.
    \item \textbf{Para diseño RF modular o multi-SDR:} usar el módulo SP4T F2914.
    \item \textbf{Para recepción solamente:} cualquiera sirve; priorizar facilidad de control.
    \item \textbf{Para transmisión o señales fuertes:} confirmar potencia máxima, aislamiento y usar cargas de 50 $\Omega$ en rutas no usadas.
\end{itemize}

\begin{advertencia}
\textbf{Nota de seguridad:} ningún switch reemplaza atenuadores, limitadores o filtros de protección. Antes de conectar un SDR, validar potencia RF, continuidad DC y estado de las rutas.
\end{advertencia}

% =========================================================
\section{Fuentes}
% =========================================================

\begin{itemize}
    \item \href{https://greatscottgadgets.com/hackrf/operacake/}{Great Scott Gadgets -- Opera Cake}.
    \item \href{https://lab401.com/products/opera-cake-for-hackrf}{Lab401 -- Opera Cake}.
    \item \href{https://www.renesas.com/en/products/rf-products/rf-switches/f2914-high-reliability-sp4t-rf-switch}{Renesas -- F2914 High Reliability SP4T RF Switch}.
    \item \href{https://www.amazon.com/Switch-Module-Single-Impedance-Reliability/dp/B0BXVVLGFB}{Amazon -- módulo RF switch SP4T}.
\end{itemize}

\end{document}
